% vim: ai sw=2 spell

\documentclass[xcolor=dvipsnames,handout]{beamer}


\let\Oldemph\emph
\renewcommand{\emph}[1]{\textcolor{blue}{\Oldemph{#1}}}
\newcommand{\heading}[1]{\textbf{\textcolor{BurntOrange}{#1}}}
\newcommand{\header}[1]{\textbf{\texttt{\textcolor{WildStrawberry}{#1}}}}
\newcommand{\doc}[1]{\textbf{\texttt{\textcolor{OliveGreen}{#1}}}}

\mode<presentation>
{
\usetheme{Madrid}
}

\setbeamertemplate{navigation symbols}{}

\usepackage[utf8x]{inputenc}
\usepackage{helvet}
\usepackage{mathptmx}
\usepackage[T1]{fontenc}
\usepackage{graphicx}
\usepackage{hyperref}
\usepackage[final]{listings}
\usepackage{multirow}
\usepackage{tikz}
\usepackage[absolute,overlay]{textpos}
\usepackage{xspace}
\usetikzlibrary{arrows, positioning}

\tikzset{tl_state/.style={rounded corners, thick, draw, minimum height=6mm,
	minimum width=12mm, bottom color=black!15, top color=black!5}}
\tikzset{tl_phase/.style={thick, draw, minimum size=8mm}}
\tikzset{tl_edge/.style={<-, >=stealth', semithick}}
\tikzstyle{na} = [baseline=-.5ex]

\lstset{tabsize=2,numbers=none,showstringspaces=false,breaklines=true,
	breakatwhitespace=true,escapeinside={(*@}{@*)},
	basicstyle=\scriptsize,keywordstyle=\bfseries\color{OliveGreen},
	commentstyle=\color{blue}, columns=flexible, language=C}

\newcommand{\cmd}[2][]{\lstinline[columns=fullflexible, language=bash,
	basicstyle=\ttfamily, #1]$#2$}
\newcommand{\code}[2][]{\lstinline[columns=fullflexible, basicstyle=\ttfamily,
	#1]$#2$}
\newcommand{\file}[2][]{\code[language=XML, #1]{#2}}

\title{Static Analysis}
\subtitle{Linux Kernel \& Used Tools}

\author{Jiri~Slaby}

\institute[SUSE Labs and FI MU]{SUSE Labs \\ and \\ Faculty of Informatics,
  Masaryk University}

\date{21. 10. \the\year}

\begin{document}

\frame{\titlepage}
\begin{frame}{Presentation Outline}
  \tableofcontents
\end{frame}

\section{Introduction}
\frame{\sectionpage}

\begin{frame}[fragile]
  \frametitle{Introduction}

  \heading{What is static analysis?}
  \begin{itemize}
    \item Analysing code similar to code inspection
    \begin{itemize}
      \item In whatever language (even machine code)
    \end{itemize}

    \pause

    \item Without directly executing code on processor
    \begin{itemize}
      \item Program state simulation
    \end{itemize}
  \end{itemize}

  \pause
  \medskip

  \begin{block}{Example}
    \begin{center}
    \begin{tabular}{c}
    \begin{lstlisting}
int foo(int a) {
  int b = 0;
  if (usually)
    b = 10;
  return a / b;
}
    \end{lstlisting}
    \end{tabular}
    \end{center}
  \end{block}

\end{frame}

\begin{frame}{Introduction}

  \heading{Why to use static analysis?}
  \begin{itemize}
    \item Can catch bugs in unlikely branches
    \begin{itemize}
      \item Usually much higher coverage
    \end{itemize}

    \pause

    \item Does not need specific hardware
    \begin{itemize}
      \item E.g.\ when checking drivers
    \end{itemize}

    \pause

    \item For debugging
    \begin{itemize}
      \item Where does this value in this variable come from?
      \item Show me only code affecting this line/variable
    \end{itemize}

  \end{itemize}

\end{frame}

\begin{frame}{Introduction}

  \heading{Why not?}
  \begin{itemize}
    \item False positives
    \begin{itemize}
      \item Imprecision in the state simulation
    \end{itemize}

    \pause

    \item May be slow
    \begin{itemize}
      \item Complex algorithms
      \begin{itemize}
	\item E.g.\ too precise state simulation to avoid false positives
      \end{itemize}

      \item Complex code
      \begin{itemize}
	\item E.g.\ find \code{y} such that \code{MD5(y) == 1} holds
      \end{itemize}
    \end{itemize}

    \pause

    \item The problem is \emph{generally undecidable}
    \begin{itemize}
      \item Not \emph{every} program can be reasoned about

      \pause

      \item This will not stop us to try our best for the rest :)
    \end{itemize}

  \end{itemize}

\end{frame}

\begin{frame}{Scope of the Presentation}

  \heading{In this presentation}
  \begin{itemize}
    \item Simple analyses
    \item Abstract interpretation
    \item Symbolic execution
  \end{itemize}

  \pause
  \bigskip

  \heading{\textbf{Not} in this presentation}
  \begin{itemize}
    \item Commercial tools
    \item Model checking
    \item \dots
  \end{itemize}

\end{frame}

\section{Techniques and Tools}
\frame{\sectionpage}

\subsection{Pattern Matching}

\begin{frame}[fragile]{Pattern Matching}

  \begin{itemize}
    \item Straightforward approach
    \item \textsc{Grep}

    \pause

    \item \textsc{Coccinelle}
    \item \textsc{FindBugs} (Java)
  \end{itemize}

  \pause
  \medskip

  \begin{block}{\textsc{Grep} Example}
  \begin{lstlisting}[language=XML]
git grep -E '%[ +0#-]*#[ +0#-]*(\*|[0-9]+)?(\.(\*|[0-9]+)?)?p'
  \end{lstlisting}
  \footnotesize{(See commit f7f94af080f2ff5cf558908c7fcc7a66f0c8adf8.)}
  \end{block}

\end{frame}

\begin{frame}[fragile]{Pattern Matching}{\textsc{Coccinelle}}

  \begin{itemize}
    \item Semantic patches
    \begin{itemize}
      \item Search ``this'' in ``this'' context
      \item And replace by ``this''

      \pause

      \item Over 50 patches ready for immediate use
      \begin{itemize}
	\item See \file{scripts/coccinelle/}
      \end{itemize}
    \end{itemize}

    \item Integrated into the Linux kernel build system
    \begin{itemize}
      \item \cmd{make coccicheck MODE=report}
      \item $\Rightarrow$ Easy to use
    \end{itemize}

    \item SmPL language

    \pause

    \item Very useful for changing API
  \end{itemize}

  \pause
  \medskip

  \begin{block}{Example}
  \begin{lstlisting}[language=XML]
$ make coccicheck MODE=report
...
drivers/staging/lustre/lustre/llite/dir.c:1055:8-15: WARNING opportunity for memdup_user
  \end{lstlisting} %stopzone
  \end{block}

\end{frame}

\subsection{Data-Flow Analysis}

\definecolor{darkblue}{rgb}{0.1, 0.1, 0.7}

\begin{frame}[fragile]
  \frametitle{Data-Flow Analysis in Example}

  \begin{center}
  \heading{Analyze possible values}

  \medskip

  \begin{tabular}{|c|c|c|c|c|} \hline
  \multirow{2}{*}{Code} & \multicolumn{2}{|c|}{\uncover<2->{must-be}} &
    \multicolumn{2}{|c|}{\uncover<3->{may-be}} \\ \cline{2-5}
  & \textcolor{darkblue}{\code{a}} & \textcolor{red}{\code{b}} &
    \textcolor{darkblue}{\code{a}} & \textcolor{red}{\code{b}} \\ \hline
  \begin{lstlisting}
void foo() {
  int (*@\textcolor{darkblue}{a}@*), (*@\textcolor{red}{b}@*) = 0;
  if (usually)
    (*@\textcolor{darkblue}{a}@*) = 1;
  else
    (*@\textcolor{darkblue}{a}@*) = 0;
  if ((*@\textcolor{red}{b}@*))
    printf("%d\n", (*@\textcolor{red}{b}@*) / (*@\textcolor{darkblue}{a}@*));  (*@\textcolor{white}{'}@*)
}
  \end{lstlisting} &
  \begin{uncoverenv}<2->
  \begin{lstlisting}
{}
{}
{}
{ 1 }
{}
{ 0 }
{}
{}
{}
  \end{lstlisting}
  \end{uncoverenv}
  &
  \begin{uncoverenv}<2->
  \begin{lstlisting}
{}
{ 0 }
{ 0 }
{ 0 }
{ 0 }
{ 0 }
{ 0 }(*@ \tikz[remember picture, na]\node[coordinate] (B0) {X}; @*)
{ 0 }
{ 0 }
  \end{lstlisting}
  \end{uncoverenv}
  &
  \begin{uncoverenv}<3->
  \begin{lstlisting}
{}
{}
{}
{ 1 }
{}
{ 0 }
{ 0, 1 }
{ 0, 1 }
{}
  \end{lstlisting}
  \end{uncoverenv}
  &
  \begin{uncoverenv}<3->
  \begin{lstlisting}
{}
{ 0 }
{ 0 }
{ 0 }
{ 0 }
{ 0 }
{ 0 }
{ 0 }
{ 0 }
  \end{lstlisting}
  \end{uncoverenv}
  \\ \hline
  \end{tabular}
  \end{center}

%  \tikz[remember picture, na]\node (X) {X};
%
%  \begin{tikzpicture}[overlay]
%    \path (X) edge (B0);
%  \end{tikzpicture}

\end{frame}

\begin{frame}{Data-Flow Analysis}

  \begin{itemize}
    \item Described by Kildall in 1970s
    \begin{itemize}
      \item \textit{A Unified Approach to Global Program Optimization}
    \end{itemize}

    \pause

    \item Intensively used for decades
    \item Optimizing compilers

    \pause

    \item Static analysis
    \begin{itemize}
      \item Sparse, Smatch (C)
      \item Clang analyzer (C, C++)
      \item Lint, Splint (C) -- dead
      \item \dots
    \end{itemize}
  \end{itemize}

\end{frame}

\begin{frame}[fragile]{Data-Flow Analysis}{\textsc{Gcc}}

  \begin{itemize}
    \item Everybody should use their favourite -Wxyz
    \item Integrated into the Linux kernel build system
    \begin{itemize}
      \item By definition
      \item $\Rightarrow$ Easy to use
    \end{itemize}
  \end{itemize}

  \pause
  \medskip

  \begin{block}{Example}
  \begin{lstlisting}[language=XML]
$ make
...
drivers/hwmon/applesmc.c: In function 'applesmc_init_smcreg':
drivers/hwmon/applesmc.c:595:43: warning: 'right_light_sensor' may be used uninitialized in this function [-Wmaybe-uninitialized]
s->num_light_sensors = left_light_sensor + right_light_sensor;
                                            ^
  \end{lstlisting} %stopzone
  \end{block}

\end{frame}

\begin{frame}[fragile]{Data-Flow Analysis}{\textsc{Sparse}}

  \begin{itemize}
    \item ``Semantic Parser''
    \item Designed for development and debugging of the Linux kernel
    \begin{itemize}
      \item Many kernel-specific checks (address spaces, etc.)
      \item Locking policy
      \item Use of \code{0} instead of \code{NULL}
      \item \dots
    \end{itemize}

    \item Integrated into the Linux kernel build system
    \begin{itemize}
      \item \cmd{make C=1}
%      \item $\Rightarrow$ Easy to use
    \end{itemize}

    \item Framework for others
  \end{itemize}

  \pause
  \medskip

  \begin{block}{Example}
  \begin{lstlisting}[language=XML]
$ make C=1
...
drivers/staging/android/ion/ion.c:1475:21: warning: incorrect type in assignment (different base types)
drivers/staging/android/ion/ion.c:1475:21:    expected restricted gfp_t [usertype] gfp_mask
drivers/staging/android/ion/ion.c:1475:21:    got int
  \end{lstlisting} %stopzone
  \end{block}

\end{frame}

\begin{frame}[fragile]{Data-Flow Analysis}{\textsc{Smatch}}

  \begin{itemize}
    \item ``Source Matcher''
    \item Based on \textsc{Sparse} framework
    \begin{itemize}
      \item Aimed at Linux kernel too
      \item More checks
      \item DMA-on-stack, wrong arguments to API functions, \dots
    \end{itemize}

    \item Integrated into the Linux kernel build system
    \begin{itemize}
      \item \cmd{make C=1 CHECK="smatch -p=kernel"}
%      \item $\Rightarrow$ Easy to use
    \end{itemize}
  \end{itemize}

  \pause
  \medskip

  \begin{block}{Use}
  \begin{lstlisting}[language=XML]
$ make C=1 CHECK="smatch -p=kernel"
...
drivers/staging/comedi/drivers/pcl812.c:738 pcl812_ai_cmd() error: we previously assumed 'dma' could be null (see line 713)
  \end{lstlisting} %stopzone
  \end{block}

\end{frame}

\begin{frame}[fragile]{Data-Flow Analysis}{\textsc{Clang}}

  \begin{itemize}
    \item From LLVM framework
    \item Besides compiler, contains also a static checker

    \pause

    \item \emph{Does not work} for Linux kernel yet
  \end{itemize}

\end{frame}

\begin{frame}[fragile]
  \frametitle{Data-Flow Analysis}
  \framesubtitle{Caveat for Static Analysis}

  \begin{center}
  \heading{How to represent multiple values?}

  \medskip

  \begin{tabular}{c}
  \begin{lstlisting}
int a = 0;
while (something) {
  a++;
}
  \end{lstlisting}
  \end{tabular}

  \begin{equation*}
    \begin{aligned}
      may\_be(a) = &\ \{ 0, 1, 2, \dots \} \\
      count(may\_be(a)) = &\ \infty
    \end{aligned}
  \end{equation*}
  \end{center}

\end{frame}

\subsection{Abstract Interpretation}

\begin{frame}{Abstract Interpretation}

  \begin{itemize}
    \item Described by Cousot\&Cousot in 1970s
    \begin{itemize}
      \item \textit{Abstract interpretation: a unified lattice model for static
	analysis of programs by construction or approximation of fixpoints}
    \end{itemize}

    \item Handle states of programs abstractly
    \item Forget irrelevant information
    \item Arbitrary levels of abstraction
  \end{itemize}

  \medskip
  \pause

  \begin{center}
  \begin{tabular}{|l|l|} \hline
  \textbf{Original code} & \textbf{Abstraction} \\ \hline
  \code{int}	& \code{Aint} \\
  \code{0}	& \code{ZERO} \\
  \code{..., -2, -1, 1, 2, ...} & \code{NONZERO} \\
  \code{int a = 5 + 0;} & \code{Aint a = NONZERO + ZERO} \\ \hline
  \end{tabular}
  \end{center}

\end{frame}

\begin{frame}[fragile]
  \frametitle{Abstract Interpretation}
  \framesubtitle{Example}

  \begin{tabular}{|l|l|l|} \hline
  Original Code & First Abstraction & Second Abstraction \\ \hline
  \begin{lstlisting}
void foo() {
  int a;
  switch (something) {
    case 0: a = 5; break;
    case 1: a = -1; break;
    case 2: a = 0; break;
  }
  if (a == 0)
    BUG();
  else if (a == -2)
    BUG();
}
  \end{lstlisting} &
  \begin{lstlisting}
void foo() {
  Aint a;
  switch (something) {
    case 0: a = POS; break;
    case 1: a = NEG; break;
    case 2: a = ZERO; break;
  }
  if (a == ZERO)
    BUG();
  else if (a == NEG)
    BUG();
}
  \end{lstlisting} &
  \begin{lstlisting}
void foo() {
  int a;
  switch (something) {
    case 0: a = POS0; break;
    case 1: a = NEG; break;
    case 2: a = 0; break;
  }
  if (a == POS0)
    BUG();
  else if (a == NEG)
    BUG();
}
  \end{lstlisting} \\ \hline
  \end{tabular}

  \bigskip
  \pause

  \begin{center}
    \heading{Be careful with abstractions!}
  \end{center}

\end{frame}

\begin{frame}{Abstract Interpretation}{\textsc{Stanse}}

  \begin{itemize}
    \item Can catch unlikely branches bugs
  \end{itemize}

\end{frame}

\subsection{Symbolic Execution}

\begin{frame}
  \frametitle{Symbolic Execution}

  \begin{itemize}
    \item Described by King in 1970s
    \begin{itemize}
      \item \textit{Symbolic execution and program testing}
    \end{itemize}
  \end{itemize}

\end{frame}

\begin{frame}{Symbolic Execution}{\textsc{Symbiotic}}

  \begin{itemize}
    \item XX
  \end{itemize}

\end{frame}

\begin{frame}{Why to use different tools?}

  \begin{itemize}
    \item Different kind of errors
    \item Different found subsets
    \item Different speed
  \end{itemize}

\end{frame}

\section{Conclusion}

\begin{frame}{Conclusion}
  XXX \\

  \bigskip
  \pause

  \heading{Questions?}
\end{frame}

\end{document}
